\documentclass[12pt, twocolumn, landscape, a4paper]{article}

\usepackage[margin=1.25cm]{geometry}
\usepackage[utf8]{inputenc}
\usepackage[spanish,es-noshorthands]{babel}
\usepackage{amssymb, amsmath, amsbsy}
\usepackage{verbatim}
\usepackage{mathpazo}
\usepackage{tasks}
\usepackage{enumerate}
\usepackage[pdftex]{hyperref}
\hypersetup{pdftitle={Anatomía Humana, CEPRU 3er examen / ciclo ordinario 2017-I, Grupo B, tema Q},pdfauthor={Luis Carrillo Gutiérrez}}

\renewcommand{\familydefault}{\sfdefault}
\pagestyle{empty} 

\begin{document}

\subsection*{Anatomía Humana}
\begin{itemize}
\item{Respecto a los \emph{órganos linfoideos secundarios}, en las siguientes proposiciones, escribir (\texttt{V}) si es verdadera o (\texttt{F}) si es falsa.
\begin{enumerate}[I.]
\item{El bazo se ubica a nivel del hipocondrio izquierdo.}
\item{El parénquima del ganglio linfático está formado por pulpa blanca y pulpa roja.}
\item{Los ganglios linfáticos son pequeños órganos de forma arriñonada.}
\item{La principal función del bazo es la filtración de la linfa por medio de la pulpa blanca.}
\end{enumerate}
La secuencia correcta es :
\begin{tasks}(5)
\task{VFVV}\task{VFVF}\task{VVFF}\task{FVFV}\task{FFVF}
\end{tasks}
}
\item{El paso del agua y solutos del plasma sanguíneo, desde los capilares glomerulares hasta el espacio de la cápsula de Bowman, se denomina :
\begin{tasks}(2)
\task{Reabsorción tubular}
\task{Secreción tubular}
\task{Filtración glomerular}
\task{Reabsorción capilar}
\task{Secreción capilar}
\end{tasks}
}
\item{En la vulva de la mujer, los pliegues longitudinales de piel cubiertos por vello púbico, homólogos del escroto del varón, son denominados :
\begin{tasks}(2)
\task{Labios mayores}
\task{Labios menores}
\task{Bulbos del vestíbulo}
\task{Cuerpos cavernosos}
\task{Cuerpos esponjosos}
\end{tasks}
}
\item{La hormona liberada por la adenohipófisis que estimula el desarrollo adicional de los folículos ováricos y la secreción máxima de los estrógenos que origina la ovulación, es :
\begin{tasks}(3)
\task{Relaxina}\task{Progesterona}\task{Melatonina}\task{Luteinizante}\task{Testosterona}
\end{tasks}
}
\end{itemize}
\end{document}